% Konfiguracja TikZ (zgodna ze skillem /tikz-diagrams) oraz pgfplots.

\usepackage{tikz}
\usetikzlibrary{
  babel,
  positioning,
  arrows.meta,
  calc,
  fit,
  backgrounds,
  shapes.geometric,
  shapes.symbols,
  shapes.multipart
}

\usepackage{pgfplots}
\pgfplotsset{compat=1.18}

% --- Paleta kolorów akademickich (/tikz-diagrams) ---
\definecolor{acadInk}{HTML}{1A1A1A}
\definecolor{acadMuted}{HTML}{595959}
\definecolor{acadFill}{HTML}{F7F7F8}
\definecolor{acadBlue}{HTML}{1F4E79}
\definecolor{acadBlueFill}{HTML}{EDF3F8}
\definecolor{acadRed}{HTML}{9B2C2C}
\definecolor{acadTeal}{HTML}{236B6E}

% \AffIcon[rozmiar]{nazwa_ikony_bez_pdf}
% Wstawia wektorową ikonę Affinity (square_gray) z katalogu images/icons/<nazwa>.pdf.
\newcommand{\AffIcon}[2][1.15cm]{%
  \IfFileExists{images/icons/#2.pdf}{%
    \includegraphics[width=#1, height=#1, keepaspectratio]{images/icons/#2.pdf}%
  }{%
    \PackageWarning{report-template}{Brak pliku ikony images/icons/#2.pdf}%
    \fbox{\parbox[c][#1][c]{#1}{\centering\tiny\texttt{#2}}}%
  }%
}

% --- Bazowe style węzłów, kart podzielnych, linii życia i przepływów (/tikz-diagrams) ---
\tikzset{
  every picture/.append style = {
    execute at begin picture = {\shorthandoff{"}}
  },
  >= {Stealth[length=2.2mm, width=1.5mm]},
  % Ostre narożniki (0pt) i białe tło (styl publikacji IEEE / USENIX / ACM)
  block/.style = {
    rectangle,
    draw=acadInk,
    semithick,
    fill=white,
    minimum width=2.4cm,
    minimum height=0.85cm,
    inner sep=5pt,
    align=center,
    font=\scriptsize
  },
  accent/.style = {
    block,
    draw=acadBlue,
    fill=acadBlueFill
  },
  % Karta z oddzielonym nagłówkiem (jak klasa UML / karta architektury)
  splitcard/.style = {
    rectangle split,
    rectangle split parts=2,
    draw=acadInk,
    semithick,
    fill=white,
    inner sep=5pt,
    text width=5.0cm,
    align=flush left,
    rectangle split part align={center, left},
    rectangle split part fill={acadFill, white},
    font=\scriptsize
  },
  splitcard_accent/.style = {
    splitcard,
    draw=acadBlue,
    rectangle split part fill={acadBlueFill, white}
  },
  % Węzeł ikony Affinity (bez ramki i bez wewnętrznego odstępu - sama grafika PDF)
  iconnode/.style = {
    inner sep=0pt,
    outer sep=2pt
  },
  % Czysty podpis pod ikoną urządzenia (bez ramki wokół tekstu)
  iconlabel/.style = {
    align=center,
    font=\scriptsize,
    inner sep=2pt,
    text=acadInk
  },
  % Nagłówek uczestnika i pionowa linia życia w diagramach sekwencji (Sequence Diagrams)
  seqactor/.style = {
    rectangle,
    draw=acadInk,
    semithick,
    fill=white,
    minimum width=3.4cm,
    minimum height=0.95cm,
    inner sep=5pt,
    align=center,
    font=\scriptsize
  },
  lifeline/.style = {
    semithick,
    densely dashed,
    draw=acadMuted
  },
  % Granica strefy bezpieczeństwa (tylko obrys, zawsze bez szarego tła)
  zone/.style = {
    rectangle,
    draw=acadMuted,
    semithick,
    densely dashed,
    fill=none,
    inner sep=8pt
  },
  % Krawędzie przepływów
  flow/.style = {
    ->,
    thick,
    draw=acadInk
  },
  ctrl/.style = {
    ->,
    semithick,
    densely dashed,
    draw=acadMuted
  },
  alert/.style = {
    ->,
    thick,
    draw=acadRed
  },
  flowlabel/.style = {
    font=\scriptsize,
    fill=white,
    inner sep=1.5pt,
    text=acadInk
  },
  step/.style = {
    circle,
    draw=acadInk,
    fill=white,
    inner sep=1pt,
    font=\tiny\bfseries,
    text=acadInk
  },
  step/.default = {}
}

% --- Domyślny styl wykresów pgfplots spójny z paletą acad* ---
\pgfplotscreateplotcyclelist{academic}{
  {acadBlue, thick, mark=*, mark size=1.6pt},
  {acadRed, thick, mark=square*, mark size=1.5pt},
  {acadTeal, thick, mark=triangle*, mark size=1.8pt},
  {acadInk, semithick, densely dashed, mark=o, mark size=1.6pt}
}

\pgfplotsset{
  every axis/.append style = {
    width = 0.85\linewidth,
    height = 6.2cm,
    axis line style = {acadInk, semithick},
    tick style = {acadInk, thin},
    label style = {font=\small, text=acadInk},
    tick label style = {font=\footnotesize, text=acadInk},
    legend style = {
      font=\footnotesize,
      draw=acadMuted,
      fill=white,
      cells={anchor=west}
    },
    grid = major,
    major grid style = {densely dashed, acadMuted!30},
    cycle list name = academic
  }
}

% \TikzFigure[szerokość]{plik_tex_w_images}{podpis}{etykieta}
% Domyślnie skaluje diagram do 0.9\linewidth; pusty argument opcjonalny [] zachowuje skalę 1:1.
\newcommand{\TikzFigure}[4][0.9]{%
  \begin{figure}
    \centering
    \IfFileExists{images/#2}{%
      \IfEmpty{#1}{%
        \input{images/#2}%
      }{%
        \resizebox{#1\linewidth}{!}{\input{images/#2}\unskip}%
      }%
    }{%
      \PackageWarning{report-template}{Brak pliku diagramu images/#2}%
      \fbox{\parbox[c][5cm][c]{0.8\linewidth}{\centering\textit{Brak pliku diagramu}\\[1ex]\code{images/#2}}}%
    }%
    \caption{#3}\label{#4}%
  \end{figure}%
}
